\section{Introduction}\label{sec:v16-introduction}
Statistical comparison asks which decisions can be transferred from
one experiment to another. When observations arrive sequentially and
the decision procedure has a fixed retained register, comparison has
an additional constraint: the randomization used to transfer an
experiment must itself obey the information and memory interface.
A hidden persistent selector, an exposed selector, and transient private
randomness define different feasible sets. Equality of unpriced
experiments does not imply equality of their fixed-register risk bodies.

We develop comparison at the original resource signature. The central
finite object is the image of a product of stochastic-row simplexes
under the map taking an executable causal simulator to its marked
feedback laws. This image is generally nonconvex. Convex separation
of its convex hull can therefore lose precisely the private-resource
obstruction one wishes to measure. Our main finite result replaces a
fixed mixture of tests by polynomial proof functions on the simulator
rows. The resulting strictly positive inequalities admit finite exact
certificates without adding a persistent selector to the simulator.

The same lower certificates can be attached to genuine physical
acquisitions. This requires more than approximating an observable in
norm: the physical approximation must preserve the actual mark,
respect feedback, produce finite reports, and rationalize their joint
laws. We retain each of those steps and prove an explicit realization
in which a nonconserved collision-compatible observable has both
nonzero generator residual and a positive exactly characterized marked
deficiency. The statistical and microscopic conclusions are joined at
the same register and selector budgets.

\subsection{The main theorem and its components}
For a fixed finite marked experiment and resource signature $b$, write
\[
 \delta_b=\min_{q\in\mathcal P_b}\max_i f_i(q).
\]
Here $q$ is an executable stochastic-row table, the $f_i$ are the
finite marked event discrepancies under deterministic feedback, and
$\mathcal P_b$ is a product of row simplexes. Visible selectors remain
part of the target reference as well as the simulated joint law.
An onto polynomial stick-breaking parameterization $q=q(x)$ places
the domain on $[0,1]^D$. With $h_i=f_i\circ q$ and $g_i=(1+h_i)/2$,
including the empty test, we prove
\begin{equation}\label{eq:v16-introdual}
 \delta_b=\sup_{p\ge1}\min_{x\in[0,1]^D}
       \frac{\sum_i g_i(x)^p h_i(x)}{\sum_i g_i(x)^p}.
\end{equation}
The minima increase to $\delta_b$ with error at most
$2\log m/(p+1)$, where $m$ is the number of tests. For rational $L$,
the strict inequality $L<\delta_b$ is equivalent to a strictly positive
finite tensor Bernstein expansion of
$\sum_i g_i^p(h_i-L)$ for some finite $p$ and degree. The proof includes
an explicit degree bound in terms of the monomial coefficients and the
positivity margin. This is Theorem~\ref{thm:v16-dual}.

The weights in \eqref{eq:v16-introdual} depend on a candidate table;
they are not independent test probabilities in an exchanged minimax.
A two-time example has private width-one deficiency $\sqrt2-1$,
hidden-two-selector deficiency zero, and visible-selector deficiency
$\sqrt2-1$ at every finite selector size. Every fixed test mixture has
minimum at most zero there. Thus the distinction is visible already
in a concrete exact problem, not just in a formal choice of dual
notation. The source supplies a strictly positive rational Bernstein
certificate for the private lower bound $2/5$.

For physical acquisitions, Theorem~\ref{thm:v16-main} combines the
polynomial certificate with stateless marked presentations. A finite
lower--upper interval $[L,U]$ gives the original physical interval
\[
 [\max\{0,L-a-c\},\ \min\{1,U+a+c\}],
\]
where $a,c$ include analytic approximation, report quantization and
joint-prefix rationalization. There is no change in the resource
signature. Under the stated effective data hypotheses these intervals
converge. The finite search and the analytic presentation search are
distinguished; neither an abstract graph core nor a decimal transition
table is called a rigorous physical input.

Our explicit microscopic realization uses $N\ge2$ interacting hard
spheres in equilibrium on a torus. For
$\vartheta=(2\pi/\ell)\sum_iq_{i,1}$ and normalized total momentum $X$,
the nonconserved observable $f=\cos\vartheta$ obeys
$U(t)f=\cos(\vartheta+\omega Xt)$ with $\omega=2\pi\sqrt N/\ell$.
All generator powers and their Gram entries are computed on the
actual collision-compatible domain. Reversal interventions give exactly
$k+1$ distinct mean functions among $2^k$ words. A Taylor presentation
then has a computable error tending to zero, uniformly over consumer
width. The actual initial sign mark yields an exact physical deficiency
series; for the parameters of Corollary~\ref{cor:v16-numerical} its
value lies strictly between $0.205809122887$ and $0.205809122889$, while
the normalized one-dimensional residual is $2\pi^2/25>0$.
The same example supplies both quantities.

Finally, for a bounded marked payoff of oscillation $M$, a variation
error $\varepsilon$ changes a normalized logarithmic value at scale
$\beta$ by at most
$\beta^{-1}\log(1+\varepsilon(e^{\beta M}-1))$.
Theorem~\ref{thm:v16-exponential} proves sharpness and transports
optimized common-policy values. This identifies the additional
accuracy needed at an exponential scale rather than asserting an
unproved large-deviation conclusion from unscaled comparison.

\subsection{Classical ingredients and the comparison of scope}
Blackwell--Le Cam comparison and its decision-theoretic formulations
provide the unpriced background \cite{Blackwell1953,LeCam1964,Torgersen1991}.
Filtered comparison is a direct antecedent, not a remote analogy.
Weisshaupt's generalized filtered randomization theorem uses a compact
convex class of filtration-preserving stochastic operators and a
separation/minimax argument \cite{Weisshaupt2002}. Our private feasible
image is not that convex class. The two-time example above states
exactly what ordinary constant convex testing loses, and the positive
polynomial theorem states exactly what replaces it. We do not claim
that polynomial positivity, power-mean approximation, or Bernstein
basis elevation is new; their elementary proofs are included to expose
the hypotheses on which the new causal application depends.

The retained existence theory credits weak-star policy topology and
strategic-measure approaches \cite{Saldi2017Topology,YukselSaldi2017}.
Common-information and finite-memory formulations are also prior
structural tools \cite{Witsenhausen1979,NayyarMahajanTeneketzis2013}.
State equivalence, incomplete-machine minimization, and weighted-automaton
methods are relevant to the retained exact task and synthesis results
\cite{Nerode1958,PaullUnger1959,Kiefer2020}. Those results are not
relabelled as new polynomial duality. The accompanying comparison
records which original sources were inspected and which proof-level
priority questions remain unverified, including the Norberg original.
The present novelty assertion is the stated resource-sensitive theorem
and its physical realization, not an assertion of an exhaustive
priority search across every form of filtered comparison.

The microscopic input is the classical almost-everywhere hard-sphere
flow \cite{GallagherSaintRaymondTexier}. The new collective calculation
uses a solvable mode of that interacting system; it is not a solution
of its kinetic limit. This distinction is mathematical: fixed-particle
$L^2$ control of noisy observations does not give a particle-uniform
BBGKY corrector on action sublevels. The program is represented as a
typed dependency graph, with theorem implications, operational adapters
and independent primary branches kept distinct.

\subsection{Organization}
We first develop intrinsic finite causal deficiency, its local and
behavior-dimension certificates, and the polynomial testing theorem.
We retain the endogenous task and stationary regenerative results, the
active graph-domain comparison, finite marked presentations, enriched
residual convergence and constructive collision data. The last sections
prove the nonconserved collective realization and the joint physical
certificate theorem. Every mathematical section of the preceding
canonical article is retained. The complete-development view additionally
contains the earlier mathematical bodies and historical introductions;
no earlier theorem is removed to create the new comparison structure.
